\documentclass[10pt,a4paper]{amsart}
\usepackage[margin=19mm]{geometry}
\usepackage{amsmath,amssymb,mathtools}
\usepackage[scr=rsfs,cal=euler]{mathalfa}
\usepackage{xfrac}
\usepackage[unicode,hidelinks,bookmarks=false]{hyperref}
\newcommand{\ie}{i.e.\ }
\newcommand{\cf}{cf.\ }
\newcommand{\resp}{respectively\ }
\newcommand{\Subsection}{\subsection}
\providecommand{\nobreakdash}{\nobreak\hbox{-}}
\setlength{\emergencystretch}{2em}
\multlinegap0pt
\usepackage{enumerate}
\usepackage{amssymb}
\usepackage{algorithm}\usepackage{algpseudocode}
\usepackage{float}
\newtheorem{assumption}{Assumption}
\newtheorem{theorem}{Theorem}
\title{\bf A Replica Exchange Markov Chain Monte Carlo Method for Disconnected Implicit Manifolds via Tubular Relaxation}
\author{Xuyuan Wang$^{\dagger}$ \and Donglin Han}
\date{Source: arXiv:2604.22055v1 (CC BY 4.0)}
\begin{document}
\maketitle

\noindent\emph{Source and licence.} \bf A Replica Exchange Markov Chain Monte Carlo Method for Disconnected Implicit Manifolds via Tubular Relaxation. Version arXiv:2604.22055v1 (CC BY 4.0), distributed by arXiv under the Creative Commons Attribution 4.0 International licence, http://creativecommons.org/licenses/by/4.0/, as recorded in the arXiv licence metadata for this exact version and shown on the abstract page https://arxiv.org/abs/2604.22055v1. Full licence notice: \texttt{licenses/arxiv-2604.22055-CC-BY-4.0.txt}.\par\smallskip

\noindent\textit{Editorial scope.} Counted unit: the article's theorem thm:jacobian\_gram (source0.tex lines 715--770). The article's complete original proof of it is delivered with it (source0.tex lines 772--931, one \texttt{proof} environment of 160 lines). No mathematics is written, repaired or re-derived: the statement and the proof are the article's own lines, in the article's own notation, and no retained line is substituted or re-flowed.  Retained uncounted as the source-local context the proof needs: lines 116--126; lines 130--133; lines 136--143; lines 405--410; lines 495--503; lines 525--537; lines 654--657.  Nothing was omitted from any retained span, so the package carries no omission record.  No retained line contains a citation, so the wrapper carries no bibliography and no bibliography entry.  No retained line contains a figure environment, an image-inclusion command or drawing code, so no image is delivered and the package declares no asset dependency.  Editorial formatting only, in addition: an AMS \texttt{amsart} preamble that restates the article's own declarations and macros used by the retained lines, namely the environments \texttt{assumption}, \texttt{theorem} on separate counters, together with the standard packages those lines need.  The retained environments print in this document as Assumption 1, Theorem 1 in order of appearance, whereas the article prints the same items with the numbering of its own full document.  The complete native source of the article as distributed in the arXiv e-print for this exact version is bundled unchanged as \texttt{sources/199066/source0.tex}.
\begin{assumption}[Constraint function]\label{ass:xi}
The constraint function $\xi$ satisfies:
\begin{enumerate}
    \item[(i)] $\xi \in C^\infty(\mathbb{R}^n, \mathbb R^m)$, i.e., it is smooth on the ambient space;
    \item[(ii)] $\nabla \xi(q):=[\nabla\xi_1(x),\nabla\xi_2(x),\cdots,\nabla\xi_m(x)]\in\mathbb R^{n\times m}$ is uniformly of full rank on $\mathcal{M}$, i.e., $\nabla \xi(q)^\top \nabla \xi(q) \succeq c I_m$ for some constant $c > 0$ and all $q \in \mathcal{M}$;
    \item[(iii)] The Hessian $\nabla^2 \xi_i(q)$ is uniformly bounded on 
    $\mathcal{M}$, i.e.\ there exists $C > 0$ such that 
    $\|\nabla^2 \xi_i(q)\| \le C$ for all $q \in \mathcal{M}$ and 
    $i = 1,\ldots,m$.
\end{enumerate}
\end{assumption}

\begin{equation}\label{eq:target_measure}
\mu(dq) 
= \frac{e^{-V(q)} \, \sigma_{\mathcal{M}}(dq)}{\int_{\mathcal{M}} e^{-V(q)} \, \sigma_{\mathcal{M}}(dq)},
\end{equation}

\begin{assumption}[Potential function]\label{ass:potential}
The potential function $V$ satisfies:
\begin{enumerate}
    \item[(i)] $V \in C^\infty(\mathbb{R}^n,\mathbb R)$, i.e., it is smooth on the ambient space;
    \item[(ii)] $V$ is bounded from below, i.e., there exists a constant $C \in \mathbb{R}$ such that $V(x) \ge C$ for all $x \in \mathbb{R}^n$;
    \item[(iii)] $V(x) \to \infty$ as $\|x\| \to \infty$, and $\int_{\mathbb{R}^n} e^{-V(x)} dx < \infty$.
\end{enumerate}
\end{assumption}

\begin{equation}\label{eq:dphi-1}
    D\Phi^{-1}(q,v)(\delta q, \delta v)
    = \bigl(I_n + E(q,v)\bigr)\delta q + \nabla\xi(q)\,\delta v,
    \qquad
    E(q,v) := \sum_{i=1}^m v_i\,\nabla^2\xi_i(q).
\end{equation}

\begin{equation}
\label{eq:tubular_relaxation}
\pi_\varepsilon(dx)
\;\propto\;
\exp\!\left(
- V(x)
- \frac{1}{\varepsilon}\|\xi(x)\|^2
\right) dx,
\end{equation}

\begin{theorem}[Weak Convergence of $\pi_\epsilon$]
\label{thm:tubular_convergence}
Let $\xi: \mathbb{R}^n \to \mathbb{R}^m$ be a constraint function satisfying Assumption \ref{ass:xi}, and let $V: \mathbb{R}^n \to \mathbb{R}$ be a potential function satisfying Assumption \ref{ass:potential}. Let $\pi_\varepsilon$ be defined by equation \eqref{eq:tubular_relaxation}. Then, as $\varepsilon \to 0$,
\begin{itemize}
    \item[(i)] There exists $C > 0$ such that
    \(
    \mathbb{E}_{\pi_\varepsilon}\left[\|v_x\|^2\right] \le C\varepsilon,
    \)
    where $v_x$ denotes the normal component of $x\in\mathcal{T}_\tau(\mathcal{M})$ under the decomposition $x = \Phi(q_x, v_x)$ with $q_x = P(x)$.
    
    \item[(ii)] The family of probability measures $\{\pi_\epsilon\}_{\epsilon>0}$ converge weakly to $\mu$ in \eqref{eq:target_measure} as $\epsilon \rightarrow 0$.
\end{itemize}
\end{theorem}

\begin{equation}\label{eq:exachange}
    S := \phi^{-1}\circ \widetilde S\circ \phi :
(q^{(k)},x^{(k)}) \longmapsto (q_k, x_k),
\end{equation}

\begin{theorem}[Jacobian determinant of $S$]
\label{thm:jacobian_gram}
Let 
\(
\mathcal M = \{q\in\mathbb R^n : \xi(q)=0\}
\)
be an implicit manifold whose constraint function satisfies Assumption \ref{ass:xi}. Define the Gram matrix
\(
G(q) := \nabla\xi(q)^\top \nabla\xi(q) \in \mathbb{R}^{m\times m}.
\)
Let $(q,x)\in\mathcal D$ be an admissible state with decomposition
\(
x = q_x + \nabla\xi(q_x)v,
\)
where
\(
\xi(q_x)=0.
\)
Let $U(q)\in\mathbb{R}^{n\times (n-m)}$ have orthonormal columns forming a basis of the tangent space $T_q\mathcal M$, satisfying
\(
U(q)^\top U(q)=I_{n-m}
\)
and
\(
\nabla\xi(q)^\top U(q)=0.
\) Then we have

\textnormal{(i)} The exact absolute value of the Jacobian determinant of the exchange map $S$ is
\begin{equation}\label{eq:jaco}
|J_S(q,x)| = \sqrt{\frac{|\det G(q)|}{|\det G(q_x)|}}
\cdot
\frac{
\left|\det\Big(I_{n-m} + \sum_{i=1}^m v_i U(q)^\top \nabla^2\xi_i(q) U(q)\Big)\right|
}{
\left|\det\Big(I_{n-m} + \sum_{i=1}^m v_i U(q_x)^\top \nabla^2\xi_i(q_x) U(q_x)\Big)\right|}.
\end{equation}

\textnormal{(ii)} Define the Gram determinant approximation
\begin{equation}\label{eq:appro}
|\widehat{J_S}(q,x)| :=
\sqrt{\frac{|\det G(q)|}{|\det G(q_x)|}}.
\end{equation}
Then there exists a constant $C>0$ such that the expected approximation error satisfies
\begin{equation}\label{eq:error}
\mathbb{E}_{\Pi}
\!\left(
\left\vert
|\widehat{J_S}(q,x)|
-
|J_S(q,x)|
\right\vert
\right)
\le C\sqrt{\epsilon}.
\end{equation}

\end{theorem}

\begin{proof}
We begin by deriving an explicit expression for the Jacobian determinant of the mapping $S$ in part (i). The exchange map can be written as the composition
\(
S = \phi^{-1} \circ \tilde S \circ \phi
\)
shown in equation \eqref{eq:exachange}, where the mappings are defined as follows
\begin{itemize}
\item $\phi(q,x) = (q,q_x,v)$ is the local coordinate mapping induced by the projection $\Phi$, with decomposition
\[
x = q_x + \nabla\xi(q_x)v,\qquad \xi(q_x)=0,
\]
\item $\tilde S(q,q_x,v) = (q_x,q,v)$ exchanges the two manifold points while leaving the normal coordinate $v$ unchanged,
\item $\phi^{-1}(q_x,q,v) = (q_x, q + \nabla\xi(q)v)$ reconstructs the original coordinates.
\end{itemize}
Applying the chain rule gives
\[
J_S(q,x)
=
J_{\phi^{-1}}(q_x,q,v)\,
J_{\tilde S}(q,q_x,v)\,
J_{\phi}(q,x).
\]
Since $\tilde S$ is a permutation of coordinates, its Jacobian determinant satisfies
\(
|\det J_{\tilde S}| = 1.
\)
Therefore,
\[
|J_S(q,x)|
=
\frac{
\left|\det J_{\phi^{-1}}(q_x,q,v)\right|
}{
\left|\det J_{\phi}(q,x)\right|
}.
\]
We now compute the Jacobian of $\phi^{-1}$. From equation \eqref{eq:dphi-1}, we have the Jacobian matrix of $\Phi^{-1}$ in block form
\[
J_{\Phi^{-1}}(q,v)
=
\begin{bmatrix}
I_n + \sum_{i=1}^m v_i \nabla^2\xi_i(q)
&
\nabla\xi(q)
\end{bmatrix},
\]
which is an $n\times n$ matrix.
To compute its determinant, we introduce an orthonormal basis adapted to the manifold. Let
\(
U(q)\in\mathbb R^{n\times (n-m)}
\)
have orthonormal columns spanning the cotangent space $T_q^*\mathcal M$, satisfying
\[
U(q)^\top U(q) = I_{n-m},\quad
\nabla\xi(q)^\top U(q) = 0.
\]
Define the orthogonal matrix
\(
C(q)
=
\begin{bmatrix}
U(q) & N(q)
\end{bmatrix},
\)
where $N(q) := \nabla\xi(q)\, G(q)^{-1/2}$ so that $C(q)$ is orthogonal and spans $\mathbb R^n$.
Since multiplication by an orthogonal matrix preserves determinants up to sign,
\[
|\det J_{\Phi^{-1}}(q,v)|
=
|\det(C(q)^\top J_{\Phi^{-1}}(q,v))|.
\]
Compute
\[
C(q)^\top J_{\Phi^{-1}}(q,v)
=
\begin{bmatrix}
U(q)^\top \left(I + \sum_{i=1}^m v_i \nabla^2\xi_i(q)\right) U(q)
&
U(q)^\top \nabla\xi(q)
\\
N(q)^\top \left(I + \sum_{i=1}^m v_i \nabla^2\xi_i(q)\right) U(q)
&
N(q)^\top \nabla\xi(q)
\end{bmatrix}.
\]
This matrix is block triangular ($U(q)^\top \nabla\xi(q)=0$), so its determinant equals the product of diagonal block determinants. Therefore,
\[
|\det J_{\Phi^{-1}}(q,v)|
=
\left|
\det\!\left(
I_{n-m}
+
\sum_{i=1}^m v_i U(q)^\top \nabla^2\xi_i(q) U(q)
\right)
\right|
\cdot
\sqrt{|\det G(q)|}.
\]
Applying this result to both $(q_x,v)$ and $(q,v)$ gives
\[
|J_S(q,x)|
=\sqrt{\frac{|\det G(q)|}{|\det G(q_x)|}}
\cdot
\frac{
\left|
\det\!\left(
I_{n-m}
+
\sum_{i=1}^m v_i U(q)^\top \nabla^2\xi_i(q) U(q)
\right)
\right|
}{
\left|
\det\!\left(
I_{n-m}
+
\sum_{i=1}^m v_i U(q_x)^\top \nabla^2\xi_i(q_x) U(q_x)
\right)
\right|
}.
\]
This establishes equation \eqref{eq:jaco} in part (i) of the theorem.

We now prove part (ii), which establishes the accuracy of the Gram determinant approximation.
Recall from part (i) that the exact Jacobian determinant of the exchange map satisfies
\begin{equation}\nonumber
|J_S(q,x)|
=
\sqrt{\frac{|\det G(q)|}{|\det G(q_x)|}}
\cdot
\frac{|\det A(q,v)|}{|\det A(q_x,v)|},
\end{equation}
where 
\begin{equation}\nonumber A(q,v) := I_{n-m} + \sum_{i=1}^m v_i\, U(q)^\top \nabla^2\xi_i(q) U(q). \end{equation}
Let $E(q,v) := \sum_{i=1}^m v_i\, U(q)^\top \nabla^2\xi_i(q) U(q)$. Under the assumption that the constraint Hessians $\nabla^2\xi_i$ are uniformly bounded on the manifold, there exists a constant $C > 0$ such that $\|E(q,v)\| \le C \|v\|$. Using the identity $\det(I + E) = 1 + \operatorname{tr}(E) + O(\|E\|^2)$, the ratio of the determinants in the normal bundle can be linearized as
\begin{equation}\nonumber
\frac{\det A(q,v)}{\det A(q_x,v)} = \frac{1 + \operatorname{tr}(E(q,v)) + O(\Vert v\Vert^2)}{1 + \operatorname{tr}(E(q_x,v)) + O(\Vert v\Vert^2)}  = 1 + \operatorname{tr}(E(q,v) - E(q_x,v)) + O(\Vert v\Vert^2).\end{equation}
Then the absolute error between the exact and approximate Jacobian is
\begin{equation}\nonumber
\left| |J_S(q,x)| - |\widehat{J_S}(q,x)| \right| = |\widehat{J_S}(q,x)| \cdot \left| \frac{\det A(q,v)}{\det A(q_x,v)} - 1 \right|.
\end{equation}
Substituting the previous linearization, we obtain
\begin{equation}\nonumber
\left| |J_S(q,x)| - |\widehat{J_S}(q,x)| \right| = |\widehat{J_S}(q,x)| \cdot \left| \sum_{i=1}^m v_i \Delta H_i(q, q_x) + O(\Vert v\Vert^2) \right|,
\end{equation}
where 
\begin{equation}\nonumber
    \Delta H_i(q, q_x) := \operatorname{tr}(U(q)^\top \nabla^2\xi_i(q) U(q)) - \operatorname{tr}(U(q_x)^\top \nabla^2\xi_i(q_x) U(q_x)).
\end{equation}
Taking the expectation with respect to $\Pi$ gives
\begin{equation}\nonumber
\mathbb{E}_{\Pi} [ | \sum_{i=1}^m v_i \Delta H_i + O(\Vert v\Vert^2) |] \le \sum_{i=1}^m \vert\Delta H_i\vert \cdot \mathbb{E}_{\pi_\varepsilon}[\Vert v_i\Vert] + O(\mathbb E_{\pi_\varepsilon}[\Vert v\Vert^2]).
\end{equation}
By Theorem \ref{thm:tubular_convergence}, we have 
\begin{equation}\nonumber
\mathbb{E}_{\Pi} \left[ \left| |J_S(q,x)| - |\widehat{J_S}(q,x)| \right| \right] = O(\sqrt{\epsilon}),
\end{equation}
completing the proof.
\end{proof}

\end{document}
